\section{Impresión de Reportes y Apoyo Administrativo}

Tu labor no concluye tras ingresar a un paciente a los equipos médicos. A menudo el paciente, y/o sus familiares, regresarán al mostrador solicitando versiones sólidas en papel de sus estudios finales, ya sea porque carecen de recursos digitales o simplemente prefieren transportarlos a otras clínicas presencialmente. 

Tu usuario Administrativo cuenta con **acceso de visualización libre a la totalidad de estudios completados de cualquier paciente en nuestra base** sin necesitar invitación previa (como en el caso de los Médicos Asociados).

\subsection{Buscar un Estudio Histórico y Extraer PDF}

\begin{enumerate}
    \item A través de tu Buscador en el sistema, escribe parte del nombre del paciente o los datos primordiales de referencia.
    \item Una vez accedes a la cronología del paciente, navega hacia abajo, ubica el estudio etiquetado como "Completado" y dale clic.
    \item Tendrás en pantalla el Dictamen del Radiólogo y la Firma certificada, exactamente igual a lo que observa un Doctor Externo desde su oficina.
    \item Al final del módulo, encontrarás iluminado el botón \textbf{"Descargar PDF"}.
    \item El archivo se descargara en la computadora del mostrador instantáneamente, y podrás mandarlo a la impresora local instalada para entregarle en fólder el informe institucional a quien lo requiera, ahorrando fricción logística a tus clientes.
\end{enumerate}

\subsection{Impresión de Tickets (\textbf{Límites Operativos})}

Para facilitar la estadía o reclamación del paciente tras admitirlo, tu sistema proveé además de reportes clínicos, boletos o "tickets" de trámite burocrático (confolio, turno, etc.) desde una de las ranuras de tu panel.

> \textbf{IMPORTANTE: RESTRICCIÓN TEMPORAL DE IMPRESIÓN}
> 
> RadiographXpress instrumenta candados lógicos antifraude. Cuando tú oprimes la opción de imprimir o descargar por primera vez un "ticket" o recibo referenciado a la ventanilla, el software inicia un cronómetro transaccional de estricta caducidad operativa en el fondo.
> 
> \textbf{Los tickets cuentan con tiempo límite para intentar obtenerse nuevamente.} Pasado un marco de escasos minutos después a la instrucción de impresión original, \textbf{la opción se bloqueará inexorablemente por el sistema}. 
>
> Esto está diseñado con el fin de evitar duplicaciones sin justificación aparente de reportes o copias falsas con horario aplazado tras salidas de la clínica. Por tanto, cerciórate siempre de que tu impresora local goza de tinta termal y papel, evadiendo tener que reportar incidencias tras bloquear un folio por descuido de límite temporal.

\begin{figure}[H]
    \centering
    \includegraphics[width=0.4\linewidth]{screenshots/ticket_print.png}
    \caption{Ejemplo del punto de control general o ticket.}
\end{figure}
